% Part: lambda-calculus
% Chapter: lambda-definability
% Section: truth-values

\documentclass[../../../include/open-logic-section]{subfiles}

\begin{document}

\olfileid{lam}{ldf}{tvr}
\olsection{ارزش‌های صدق و روابط}

می‌توانیم ارزش‌های صدق را در حساب لامبدای محض چنین کدگذاری کنیم:
\begin{align*}
  \fn{true} & \ident \lambd[x][\lambd[y][x]]\\
  \fn{false} & \ident \lambd[x][\lambd[y][y]]
\end{align*}

ارزش‌های صدق به‌صورت \emph{گزینشگر} نمایش داده می‌شوند؛ یعنی توابعی که
دو آرگومان می‌پذیرند و یکی از آن‌ها را بازمی‌گردانند. ارزش صدق
$\fn{true}$ آرگومان نخست و~$\fn{false}$ آرگومان دوم را برمی‌گزیند. برای
مثال، $\fn{true}\, M N$ همواره به~$M$ و~$\fn{false}\, M N$ همواره
به~$N$ کاهش می‌یابد.

\begin{defn}
رابطهٔ~$R \subseteq \Nat^k$ را~!!{lambda definable} می‌نامیم اگر ترمی
چون~$R$ وجود داشته باشد که
\begin{align*}
  R\, \num{n_1} \dots \num{n_k} & \bred \fn{true}
  \intertext{هرگاه~$R(n_1, \dots, n_k)$ برقرار باشد؛ و}
  R\, \num{n_1} \dots \num{n_k} & \bred \fn{false}
\end{align*}
در غیر این صورت.
\end{defn}

برای نمونه، رابطهٔ~$\fn{IsZero} = \{0\}$ که فقط برای~$0$ برقرار است،
به‌وسیلهٔ ترم زیر~!!{lambda definable} است:
\[
  \fn{IsZero} \ident \lambd[n][n (\lambd[x][\fn{false}])\, \fn{true}].
\]
این ترم چگونه کار می‌کند؟ چون عددنماهای چرچ به‌صورت تکرارگر تعریف
شده‌اند (توابعی که آرگومان نخست خود را~$n$ بار بر آرگومان دوم اعمال
می‌کنند)، مقدار آغازین را~$\fn{true}$ می‌گذاریم و در هر گام تکرار، بدون
توجه به نتیجهٔ تکرار قبلی، $\fn{false}$ را بازمی‌گردانیم. این گام
$n$ بار بر مقدار آغازین اعمال می‌شود؛ پس نتیجه اگر و تنها اگر هیچ گامی
اعمال نشود، یعنی وقتی~$n = 0$، برابر~$\fn{true}$ خواهد بود.

بر پایهٔ این نمایش ارزش‌های صدق، می‌توانیم چند تابع صدق را نیز تعریف
کنیم. در اینجا دو نمونه، یعنی بازنمایی‌های نقیض و عطف، آمده‌اند:
\begin{align*}
  \fn{Not} & \ident \lambd[x][x\, \fn{false}\, \fn{true}]\\
  \fn{And} & \ident \lambd[x][\lambd[y][xy \,\fn{false}]]
\end{align*}
تابع «$\fn{Not}$» یک آرگومان می‌پذیرد و اگر آرگومان~$\fn{false}$ باشد
$\fn{true}$، و اگر آرگومان~$\fn{true}$ باشد~$\fn{false}$ را
بازمی‌گرداند. تابع «$\fn{And}$» دو ارزش صدق را به‌عنوان آرگومان می‌پذیرد
و باید اگر و تنها اگر هر دو آرگومان~$\fn{true}$ باشند، $\fn{true}$ را
بازگرداند. ارزش‌های صدق به‌صورت گزینشگر نمایش داده شده‌اند؛ بنابراین
وقتی~$x$ یک ارزش صدق باشد و بر دو آرگومان اعمال شود، اگر~$x$ برابر
$\fn{true}$ باشد نتیجه آرگومان نخست، و در غیر این صورت آرگومان دوم است.
اکنون~$\fn{And}$ دو آرگومان~$x$ و~$y$ را می‌گیرد و~$y$ و~$\fn{false}$
را به آرگومان نخست خود، یعنی~$x$، می‌دهد. با فرض اینکه~$x$ یک ارزش صدق
باشد، اگر~$x$ برابر~$\fn{true}$ باشد حاصل به~$y$، و اگر~$x$ برابر
$\fn{false}$ باشد به~$\fn{false}$ ارزیابی می‌شود؛ دقیقاً همان نتیجه‌ای
که می‌خواهیم.

توجه کنید که در اینجا فرض می‌کنیم فقط ارزش‌های صدق به‌عنوان آرگومان به
$\fn{And}$ داده می‌شوند. اگر ترم‌های دیگری به آن داده شوند، حاصل، یعنی
صورت نرمال اگر وجود داشته باشد، ممکن است ارزش صدق نباشد.

\begin{prob}
  با استفاده از کدگذاری ارزش‌های صدق به‌صورت ترم‌های~$\lambd$، توابع
  $\fn{Or}$ و~$\fn{Xor}$ را تعریف کنید که به‌ترتیب توابع صدق فصل و فصل
  مانعةالجمع را نمایش می‌دهند.
\end{prob}

\end{document}
